% Endpoint-24 Affine-Defect Boundary Atlas --- paper source
%
% Public-facing claims are controlled by paper/CLAIM_BOUNDARY.md.
% Finite replay statements must point to paper/SOURCE_MAP.md rows.

\documentclass[11pt,a4paper]{article}

\usepackage[T1]{fontenc}
\usepackage[utf8]{inputenc}
\usepackage{lmodern}
\usepackage{amsmath,amssymb,amsthm,mathtools}
\usepackage{booktabs}
\usepackage{array}
\usepackage{enumitem}
\usepackage{etoolbox}
\usepackage{microtype}
\usepackage[margin=1in]{geometry}
\usepackage{xcolor}
\usepackage{xurl}
\usepackage{xspace}
\usepackage{float}
\usepackage{placeins}
\usepackage{hyperref}

\hypersetup{
  colorlinks=true,
  linkcolor=blue!55!black,
  citecolor=blue!55!black,
  urlcolor=blue!55!black,
  pdftitle={Endpoint-24 Affine-Defect Boundary Atlas},
  pdfauthor={Oleksiy Babanskyy},
  pdfsubject={Endpoint-24 affine-defect boundary atlas for order-four magic squares},
  pdfkeywords={magic squares, Durer square, F2-to-the-4, affine geometry, finite certificates}
}
\emergencystretch=2em
\setlength{\textfloatsep}{8pt plus 2pt minus 2pt}
\setlength{\floatsep}{8pt plus 2pt minus 2pt}
\setlength{\intextsep}{8pt plus 2pt minus 2pt}
\setlength{\abovecaptionskip}{3pt}
\setlength{\belowcaptionskip}{0pt}
\AtBeginEnvironment{thebibliography}{%
  \footnotesize
  \setlength{\itemsep}{0pt}%
  \setlength{\parsep}{0pt}%
  \setlength{\parskip}{0pt}%
  \setlength{\topsep}{0pt}%
}

\newtheorem{theorem}{Theorem}[section]
\newtheorem{proposition}[theorem]{Proposition}
\newtheorem{certprop}[theorem]{Certified Proposition}
\newtheorem{lemma}[theorem]{Lemma}
\newtheorem{corollary}[theorem]{Corollary}
\theoremstyle{definition}
\newtheorem{definition}[theorem]{Definition}
\theoremstyle{remark}
\newtheorem{remark}[theorem]{Remark}

\newcommand{\F}{\mathbb F}
\newcommand{\A}{\mathcal A}
\newcommand{\Derr}{\mathcal D}
\newcommand{\supp}{\operatorname{supp}}
\newcommand{\TK}{T/K}
\newcommand{\Durer}{D\"urer\xspace}

\title{Endpoint-24 Affine-Defect Boundary Atlas:\\
B1/B2/B3, \(F_2^4\) Error Geometry, and Finite Transport}
\author{Oleksiy Babanskyy}
\date{}

\begin{document}
\maketitle

\begin{abstract}
We extract a finite theorem within the certified endpoint-24
order-four magic-square atlas.  In value-minus-one coordinates over \(\F_2^4\), the affine
interpolation error \(E=L+A\) has defect \(\delta\in\{0,4,6\}\).  On endpoint
\(24\), this defect and the selected-label xor split the certified atlas into
an exact affine core, a selected-affine B2 boundary, and a selected-non-affine
B3 boundary:
\[
  236=144+32+60.
\]
The verified free quotient \(\TK\cong V_4\) gives
\[
  59=36+8+15.
\]
The B3 support-action matrix has \(64\) candidate signature/profile cells,
with exactly \(15\) realized and \(49\) absent cells.  The eight B3 profile
columns are now certified intrinsic to the selected signature and bilinear error
matrix; the remaining realization cut is certified finite and integer-arithmetic,
not an \(\F_2\)-linear closure.  We separate the first-principles part of the
argument, namely the degree/defect and error-fiber parity spine, from the finite
transport and matrix certificates.  The symbolic derivation of the B3 count
\(60\) remains an explicit open debt.
\end{abstract}

\noindent\textbf{Keywords.} Magic squares, \Durer square, affine geometry over
\(F_2\), finite certificates, endpoint atlas.

\section{Introduction}

This paper is a focused companion to the public order-four magic-square work.
Its object is narrower: the certified endpoint-24 boundary atlas in
affine-defect coordinates.  The earlier paper isolates \(24\) as a structured
bounded meet.  Here we examine the finite boundary that appears when the
order-four endpoint-24 atlas is organized by affine error over \(\F_2^4\).

The central split is
\[
  236 = 144 + 32 + 60,
\]
where the three terms are the exact affine core, the selected-affine B2
boundary, and the selected-non-affine B3 boundary.  The verified free
\(\TK\cong V_4\) action gives the quotient split
\[
  59 = 36 + 8 + 15.
\]

The exact affine core has \(144\) records.  The selected-affine boundary B2 has
\(32\) records.  The selected-non-affine boundary B3 has \(60\) records and
carries a support-action matrix with \(15\) realized and \(49\) absent
signature/profile cells.  B1 is included as a scoped companion statement: the
endpoint \(23/27\) signatures are absent from the affine selected-plane layer by
the same four-point xor criterion, without claiming a full endpoint \(23/27\)
classification.

\paragraph{Contribution.}
The paper has three layers.
\begin{enumerate}[leftmargin=*]
\item It states the endpoint-24 affine-defect boundary theorem in a compact
  certified-atlas form.
\item It separates the structural part of the proof from the finite matrix
  part.  The degree/defect dichotomy and the error-fiber parity bridge are
  proved as first-principles finite theorems; the B3 support-action matrix is
  a finite transport certificate.
\item It collapses the remaining B3 count debt from ``why \(15\) orbits'' to
  the sharper finite kernel
  \[
    \delta=0,4,6\mapsto 36,16,7
  \]
  and the B3 row-degree pattern \(1,1,1,2,2,1,4,3\).
\end{enumerate}

\paragraph{Claim discipline.}
Replay-backed statements are labelled as Certified Propositions and point to
the source map in \texttt{paper/SOURCE\_MAP.md}.  The manuscript does not claim
a non-enumerative derivation of the B3 count \(60\), a symbolic derivation of
the endpoint kernel \(36,16,7\), a full endpoint \(23/27\) classification, or a
full classification of non-affine order-four magic-square behavior.
\section{Model and Notation}

Cells are indexed by \(\F_2^4\), and square values are converted to labels by
subtracting one.  Thus the values \(1,\ldots,16\) become labels
\(0,\ldots,15\).  For a normal order-four magic square, let
\[
  L:\F_2^4\to\F_2^4
\]
be the label map.  Let \(A\) be the affine interpolation of \(L\) on the five
base points \(\{0,e_1,e_2,e_3,e_4\}\), and define the error map
\[
  E=L+A,\qquad \delta=|\supp(E)|.
\]

For four distinct selected labels, the selected set is an affine plane exactly
when its xor is zero.  This selected-xor axis is the finite selector used to
separate the exact affine row, B2, and B3 on endpoint \(24\).

\begin{definition}[Square-mask pair and endpoint]
A square-mask pair is a normal order-four magic square together with a
one-incidence mask.  The mask selects one cell in each row, one cell in each
column, and one cell on each of the two main diagonals.  The endpoint is
\(34-t\), where \(t=\min_{c\in M} L(c)\) is the minimum selected label (the
minimum selected value minus one).  Equivalently, endpoint \(24\) is exactly
the slice whose minimum selected value is \(11\).  This paper works inside the
certified endpoint \(24\) subatlas, which contains \(236\) such records.  In
particular every endpoint-\(24\) selected four-set contains the value \(11\),
which is why the B3 signatures are four-subsets of \(\{11,\ldots,16\}\)
containing \(11\).
\end{definition}

\begin{definition}[Selected mask and selected labels]
For a square-mask pair, \(M\subset \F_2^4\) denotes the four selected cells.
The selected labels are the four elements \(L(M)\subset \F_2^4\).  The selected
set is affine if and only if
\[
  \bigoplus_{c\in M}L(c)=0.
\]
Equivalently, in value notation, the selected values form an affine four-set
after subtracting one from each value.
\end{definition}

\begin{definition}[Endpoint-24 record]
An endpoint-24 record is a certified square-mask pair whose endpoint is \(24\).
The development atlas contains \(236\) such records.
\end{definition}

\begin{definition}[Cell-value affineness and essential representatives]
A record is cell-value affine when its full label map \(L:\F_2^4\to\F_2^4\) is
affine.  The finite atlas is stored modulo the usual dihedral \(D_4\) symmetries
of the square; theorem statements that count quotient objects explicitly say
when the free \(\TK\) quotient is also being used.
\end{definition}

\begin{definition}[The \(\TK\) quotient]
The finite transport group \(T\) is the row/column relation group used by the
extension certificates, and \(K\) is its loop stabilizer.  The verified quotient
\(\TK\cong V_4\) acts freely on every certified square-mask pair and preserves
the endpoint, affine-defect class, and selected signature.  Hence each
\(\TK\)-orbit has four records.
\end{definition}

\begin{definition}[Signature and support-action profile]
For B3, a signature is one of the eight non-affine selected-label four-sets
containing the minimum selected value \(11\).  A support-action profile records
the defect \(\delta\), the error-direction space \(\Derr\), and the finite
defect-mask support type seen by the selected mask.  The B3 support-action
matrix has these eight signatures as rows and the eight realized profiles as
columns.
\end{definition}

\begin{definition}[Clause templates]
The public names NF1, NF2, and NF3 refer to the three B3 support-action clause
templates.  NF1 is the defect-\(4\) singleton xor-error case.  NF2 is the
defect-\(6\) non-xor pair-support case.  NF3 is the defect-\(6\) non-xor
triple-support case.  The templates are structural; the thirteen profile-level
instances are finite-certified.
\end{definition}

\begin{definition}[Affine defect]
The affine defect of a record is \(\delta=|\supp(E)|\), where \(E=L+A\).  The
associated error direction space \(\Derr\) is the span of the bilinear error
coefficients.  The certified dichotomy is
\[
  (\delta,\dim \Derr)\in\{(0,0),(4,1),(6,2)\}.
\]
\end{definition}

\begin{definition}[Boundary labels]
Inside endpoint \(24\), we use the following names.
\begin{center}
\begin{tabular}{llrr}
\toprule
Name & Selector & Records & \(\TK\)-orbits\\
\midrule
exact affine core & selected affine, \(\delta=0\) & \(144\) & \(36\)\\
B2 & selected affine, \(\delta=4\) & \(32\) & \(8\)\\
B3 & selected non-affine, \(\delta\in\{4,6\}\) & \(60\) & \(15\)\\
\bottomrule
\end{tabular}
\end{center}
\end{definition}

The classical context is the order-four magic-square structure of Ollerenshaw
and Bondi~\cite{ollerenshawBondi}, vector-space viewpoints on magic
squares~\cite{wardMagicVectorSpaces}, the \Durer tesseract
interpretation~\cite{sudberyTesseract}, and enumeration work on magic squares
and the order-four case~\cite{beckCohenCuomoGribelyukMagic,
beckVanHerick4x4Enumeration}.  A recent
computational classification of the 880 essentially different normal
fourth-order magic squares is recorded by Takemura~\cite{takemura880}.  The
operative finite proofs in this paper are the local certificates listed in
Section~\ref{sec:certificates}.
\section{Algebraic Spine and Boundary Theorem}
\label{sec:theorem}

The theorem below states the finite boundary result in public notation.
Each finite assertion is tied to a source artifact in
\(\texttt{paper/SOURCE\_MAP.md}\); the proof uses those artifacts as certified
finite inputs rather than as hidden mathematical assumptions.

\begin{certprop}[Line-xor and degree-bound certificate]
\label{cert:line-xor-degree}
In value-minus-one labels, every row and every column of each certified normal
order-four magic square has xor zero.  This is the local Nimm0/Lemma B input.
The manuscript treats the Step4 finite derivation as the proof source; the
classical order-four literature supplies context rather than an unverified
quoted theorem~\cite{ollerenshawBondi,wardMagicVectorSpaces}.  With the
row/column line-xor input, the degree-bound argument gives
\[
  \deg L\le 2.
\]
\end{certprop}

\begin{proof}
Write each output bit of \(L(r_1,r_0,c_1,c_0)\) in algebraic normal form over
\(\F_2\).  For a monomial \(r^{S_r}c^{S_c}\), xoring over a fixed row leaves it
alive exactly when \(S_c=\{c_1,c_0\}\).  Hence the condition that every row xors
to zero kills precisely the monomials divisible by \(c_1c_0\).  Similarly,
column-xor zero kills precisely the monomials divisible by \(r_1r_0\).  Every
degree-three or degree-four monomial in the four variables is divisible by one
of these two pairs, so no monomial of degree at least three remains.  The
row/column line-xor premise itself is the Step4 finite row-pair derivation from
the full row, column, and diagonal magic constraints.
\end{proof}

\begin{lemma}[Quadratic error form]
\label{lem:quadratic-error-form}
Let \(A\) be the affine interpolation of \(L\) on
\(\{0,e_1,e_2,e_3,e_4\}\), and set \(E=L+A\).  For the certified normal
magic-square records, the pure row-pair and pure column-pair quadratic
coefficients vanish.  Hence the error has bilinear form
\[
  E(r,c)=B_{11}r_1c_1+B_{10}r_1c_0+B_{01}r_0c_1+B_{00}r_0c_0.
\]
\end{lemma}

\begin{proof}
The affine interpolation \(A\) absorbs the constant and linear parts of \(L\).
By Proposition~\ref{cert:line-xor-degree}, no terms of degree at least three
occur.  The row-xor condition kills the pure column-pair term \(c_1c_0\), while
the column-xor condition kills the pure row-pair term \(r_1r_0\).  The only
quadratic terms left are the four mixed row/column monomials displayed above.
\end{proof}

\begin{theorem}[Defect dichotomy / Lemma C]
\label{thm:defect-dichotomy}
For the degree-at-most-two bijections of \(\F_2^4\) covered by the certificate,
the bilinear coefficient matrix has the certified one-dependent-side property.
Consequently the error-direction span \(\Derr\) has dimension at most two and
\[
  (\delta,\dim\Derr)\in\{(0,0),(4,1),(6,2)\}.
\]
This is the first-principles structural part of the affine-defect spine.
\end{theorem}

\begin{proof}
Let \(B=(B_{ij})\) be the \(2\times2\) matrix of mixed bilinear coefficients.
The property that a bilinear matrix is completable to a degree-at-most-two
bijection \(x\mapsto A_0x+E(x)\), and the property that one side of \(B\) is
dependent, are invariant under the codomain action of \(GL(4,2)\).  The
Step5+ check therefore reduces the obstruction to a symmetry-complete finite
set: all \(4096\) matrices with entries in a fixed three-dimensional subspace,
plus one representative for the four-dimensional span case.  In that reduced
universe every completable matrix has one dependent side, and the
four-dimensional representative is not completable.  Thus a quadratic
bijection has \(\dim\Derr\le2\) and a dependent side.

With one dependent side, \(E\) factors as a nonzero scalar linear form on one
\(\F_2^2\) coordinate side times a vector-valued linear map on the other side.
If \(d=\dim\Derr\), the scalar form is nonzero on two points and the vector map
is nonzero on \(4-2^{2-d}\) points.  Hence
\[
  \delta=2(4-2^{2-d}),
\]
which gives exactly \((0,0),(4,1),(6,2)\).
\end{proof}

\begin{lemma}[Error-fiber parity bridge]
\label{lem:error-fiber-parity}
Let \(M\) be the selected mask.  If \(\delta=4\), then
\(\Derr=\{0,u\}\) and \(E=u\) on a four-cell support \(\Sigma\), so
\[
  \bigoplus_{c\in M} L(c)=u\cdot(|M\cap\Sigma|\bmod 2).
\]
If \(\delta=6\), the three nonzero fibers
\(\Sigma_u,\Sigma_v,\Sigma_{u+v}\) contribute
\[
  \bigoplus_{c\in M}L(c)=
  (|M\cap\Sigma_u|\bmod2)u+
  (|M\cap\Sigma_v|\bmod2)v+
  (|M\cap\Sigma_{u+v}|\bmod2)(u+v).
\]
Thus selected-label affineness is controlled by error-fiber parity.
\end{lemma}

\begin{proof}
The selected masks in this endpoint atlas are affine planes in the cell space,
so the affine part cancels:
\[
  \bigoplus_{c\in M}A(c)=0.
\]
Thus \(\bigoplus_{c\in M}L(c)=\bigoplus_{c\in M}E(c)\).  The factorization in
Theorem~\ref{thm:defect-dichotomy} gives the fiber shapes: for \(\delta=4\) a
single nonzero error value \(u\) occurs on a four-cell support, and for
\(\delta=6\) the three nonzero values \(u,v,u+v\) occur on three two-cell
fibers.  Grouping the xor of \(E\) over \(M\) by these fibers gives exactly the
two displayed parity formulas.
\end{proof}

\begin{certprop}[Free \(\TK\)-transport]
\label{cert:tk-transport}
The finite transport certificate proves a free global \(\TK\cong V_4\) action on
the certified atlas.  It preserves endpoint, class, defect, and selected
signature.  Therefore endpoint \(24\) has \(236=59\cdot4\), with quotient split
\[
  59=36+8+15.
\]
\end{certprop}

\begin{certprop}[B2 selected-affine boundary]
\label{cert:b2-boundary}
Inside endpoint \(24\), B2 is the selected-affine, cell-value-non-affine
boundary.  It has \(32\) records and \(8\) free \(\TK\)-orbits.  In the
\(\delta=4\) fiber picture it is the even side of the constant-error support
parity:
\[
  |M\cap\Sigma|\equiv0\pmod2.
\]
The endpoint package verifies this selected-affine, cell-value-non-affine
finite selector and closes the \(32/8\) count inside the saved endpoint atlas.
This is a structural parity selector plus a finite endpoint package, not a
complete conceptual B2 classification theorem.
\end{certprop}

\begin{certprop}[B3 support-action matrix]
\label{cert:b3-support-action}
Inside endpoint \(24\), B3 is the selected-non-affine boundary.  It has
\(60\) records and \(15\) free \(\TK\)-orbits.  The support-action matrix has
eight non-affine selected signatures against eight support-action profiles,
hence \(64\) candidate cells.  The finite bilinear transport package verifies
exactly \(15\) realized cells and \(49\) absent cells, with no false positives
or false negatives.  The intrinsic column-classification certificate refines
the column side: all eight profile columns are intrinsic to the selected
signature and the bilinear error matrix \(B\), through \(\delta\), the zero
pattern of \(B\), and \(\Derr=\operatorname{span}(B)\).  The defect-6
realizability-obstruction certificate identifies the remaining realization side
as an integer-magic residue: at defect \(6\), the order-four magic-square
constraints realize \(13\) of the \(35\) two-dimensional error spaces, although
all \(35\) are \(\F_2\)-completable and the realized \(13\)-set has trivial
\(GL(4,2)\) stabilizer.  Thus the column classification is intrinsic, while
the \(15/49\) realization cut remains a finite arithmetic certificate.  The
NF1/NF2/NF3 templates are structural clause shapes; the thirteen profile-level
instances are finite-certified.
\end{certprop}

\begin{certprop}[B3 count collapse]
\label{cert:b3-count-collapse}
The quotient count \(15\) is not an independent count once freeness is known:
\[
  15=60/4.
\]
The finite residual is the row-degree pattern
\[
  1,1,1,2,2,1,4,3
\]
and the endpoint quotient kernel
\[
  \delta=0,4,6\mapsto36,16,7.
\]
This is a finite debt reduction, not a symbolic proof of the count \(60\).  A
non-enumerative derivation of that count remains outside the claim.
\end{certprop}

\begin{certprop}[B1 scoped affine-plane absence]
\label{cert:b1-absence}
B1 is the scoped endpoint \(23/27\) absence statement for the affine
selected-plane layer.  In that layer every selected four-set has selected-label
xor zero.  The endpoint \(23/27\) signatures under discussion have nonzero
selected-label xor, and therefore cannot occur in the affine selected-plane
layer.  This is not a full endpoint \(23/27\) classification.
\end{certprop}

\begin{theorem}[Affine-defect boundary atlas, certified endpoint-24 form]
\label{thm:b123-affine-defect-boundary}
Inside the certified endpoint-24 order-four magic-square atlas, consider square-mask pairs in
the bounded one-incidence model.  Quotient statements are taken modulo the
verified \(D_4\) essential representatives and the verified free \(\TK\)
action.  Then:
\begin{enumerate}[leftmargin=*]
\item endpoint \(24\) has \(236\) records and \(59\) free \(\TK\)-orbits;
\item the affine defect satisfies \(\delta\in\{0,4,6\}\), with record split
  \(144/64/28\);
\item the exact affine core has \(144\) records and \(36\) \(\TK\)-orbits;
\item B2 has \(32\) records and \(8\) \(\TK\)-orbits;
\item B3 has \(60\) records and \(15\) \(\TK\)-orbits.  Its support-action
  matrix has \(64\) candidate cells, exactly \(15\) realized cells, and
  exactly \(49\) absent cells;
\item B1 is the scoped endpoint \(23/27\) absence statement from the affine
  selected-plane layer, not a full endpoint \(23/27\) classification.
\end{enumerate}
\end{theorem}

\begin{proof}
Proposition~\ref{cert:line-xor-degree} gives \(\deg L\le2\) from row/column
line xor zero, and Lemma~\ref{lem:quadratic-error-form} writes \(E=L+A\) as a
bilinear error on row/column coordinates.  Theorem~\ref{thm:defect-dichotomy}
gives \(\delta\in\{0,4,6\}\) with
\((\delta,\dim\Derr)=(0,0),(4,1),(6,2)\), and the error-fiber parity bridge
identifies the selected-label xor from the selected mask's intersections with
the error fibers.

The free \(\TK\)-transport proposition gives \(236=59\cdot4\) on endpoint
\(24\).  The exact affine core is the selected-affine, \(\delta=0\) cell and
has \(144\) records, hence \(36\) free \(\TK\)-orbits.  The B2 proposition gives
the selected-affine boundary as \(32\) records and \(8\) orbits.  The B3
support-action proposition gives the selected-non-affine boundary as \(60\)
records and \(15\) orbits, with the \(15/49\) support-action matrix.  The
count-collapse proposition records the sharper finite residual without turning
it into a symbolic count proof.  Finally, the B1 proposition gives the scoped
endpoint \(23/27\) absence from the affine selected-plane layer.  These blocks
are exactly the assertions of the theorem.
\end{proof}

\begin{certprop}[Full endpoint aggregate controls]
\label{cert:full-endpoint-control}
The public aggregate control artifact verifies the full
\(880\cdot 8=7040\) square--mask universe used as context for this paper.  For
every square--mask pair with label map \(L\) and selected mask \(M\),
\[
  \operatorname{end}(L,M)=34-\min_{c\in M}L(c)
  =35-\min\{\text{selected values}\}.
\]
The endpoint distribution is
\[
\begin{array}{c|rrrrrrrrrrrrr}
 e&22&23&24&25&26&27&28&29&30&31&32&33&34\\\hline
 \#&176&44&236&268&676&108&328&376&772&420&884&992&1760.
\end{array}
\]
The same artifact gives the aggregate \(T/K\) endpoint-orbit counts
\[
44,11,59,67,169,27,82,94,193,105,221,248,440
\]
for endpoints \(22,23,\ldots,34\).  Endpoints \(23\) and \(27\) have no
possible selected-affine signature, while endpoint \(22\) is a positive
selected-affine control: all its \(176\) aggregate records are selected-affine.
Endpoint \(24\) is therefore the first anchor treated here where the
selected-affine and selected-non-affine layers both occur, with the certified
split \(176+60=236\) and affine-defect distribution \(144,64,28\).

This is a finite aggregate control statement.  It is not a classification of
every endpoint and it does not widen
Theorem~\ref{thm:b123-affine-defect-boundary} beyond endpoint \(24\).
\end{certprop}

% Generated by scripts/generate_paper_source_backed_tables.py --write.
% Do not edit by hand; edit the generator or source JSON artifacts.

\section{Source-Backed Tables and Figures}
\label{sec:source-backed-tables}

The following compact displays are generated from the certified JSON
artifacts listed in \texttt{paper/SOURCE\_MAP.md}.  They expose only
paper-facing aggregate data, not record-level payloads.

\begin{figure}[ht]
\centering
\[
  236=144+32+60,\qquad
  59=36+8+15.
\]
\caption{Endpoint-24 record and quotient splits.  Source-map rows: endpoint split, free transport, and count collapse.}
\end{figure}

\begin{table}[ht]
\centering
\small
\begin{tabular}{llrr}
\toprule
Selected affine & Defect $\delta$ & Records & Class\\
\midrule
yes & 0 & 144 & exact affine\\
yes & 4 & 32 & B2\\
yes & 6 & 0 & empty\\
no & 0 & 0 & empty\\
no & 4 & 32 & B3\\
no & 6 & 28 & B3\\
\bottomrule
\end{tabular}
\caption{Endpoint-24 two-invariant table: selected-label affineness against affine defect.}
\end{table}

\begin{table}[ht]
\centering
\small
\begin{tabular}{lrrr}
\toprule
$\delta$ & Records & $\TK$-orbits & Preserved diagonal planes\\
\midrule
0 & 144 & 36 & 24\\
4 & 64 & 16 & 8\\
6 & 28 & 7 & 0\\
\bottomrule
\end{tabular}
\caption{Endpoint quotient kernel and finite diagonal-plane avatar.  The last column is certified finite data, not a first-principles derivation.}
\end{table}

\begin{table}[ht]
\centering
\scriptsize
\begin{tabular}{lrrrrl}
\toprule
Signature & xor & records & orbits & degree & $\delta$ split\\
\midrule
\texttt{11,12,13,15} & 3 & 4 & 1 & 1 & \texttt{d4:4}\\
\texttt{11,12,13,16} & 2 & 4 & 1 & 1 & \texttt{d4:4}\\
\texttt{11,12,14,15} & 2 & 4 & 1 & 1 & \texttt{d4:4}\\
\texttt{11,12,14,16} & 3 & 8 & 2 & 2 & \texttt{d4:4, d6:4}\\
\texttt{11,13,14,15} & 5 & 8 & 2 & 2 & \texttt{d4:4, d6:4}\\
\texttt{11,13,14,16} & 4 & 4 & 1 & 1 & \texttt{d4:4}\\
\texttt{11,13,15,16} & 7 & 16 & 4 & 4 & \texttt{d6:16}\\
\texttt{11,14,15,16} & 6 & 12 & 3 & 3 & \texttt{d4:8, d6:4}\\
\bottomrule
\end{tabular}
\caption{B3 selected-signature degrees.  The degree column is the row sum in the support-action matrix.}
\end{table}

\begin{table}[ht]
\centering
\scriptsize
\begin{tabular}{lccccccccr}
\toprule
Signature & \texttt{DM01} & \texttt{DM02} & \texttt{DM03} & \texttt{DM04} & \texttt{DM05} & \texttt{DM06} & \texttt{DM07} & \texttt{DM08} & degree\\
\midrule
\texttt{11,12,13,15} & 1 & . & . & . & . & . & . & . & 1\\
\texttt{11,12,13,16} & . & 1 & . & . & . & . & . & . & 1\\
\texttt{11,12,14,15} & . & 1 & . & . & . & . & . & . & 1\\
\texttt{11,12,14,16} & 1 & . & 1 & . & . & . & . & . & 2\\
\texttt{11,13,14,15} & 1 & . & . & . & . & . & 1 & . & 2\\
\texttt{11,13,14,16} & . & 1 & . & . & . & . & . & . & 1\\
\texttt{11,13,15,16} & . & . & . & 1 & 1 & 1 & . & 1 & 4\\
\texttt{11,14,15,16} & 1 & 1 & . & . & . & . & 1 & . & 3\\
\bottomrule
\end{tabular}
\caption{B3 support-action matrix.  Entry $1$ means realized; a dot means absent.}
\end{table}

\begin{table}[ht]
\centering
\scriptsize
\begin{tabular}{lrrll}
\toprule
Profile & $\delta$ & orbits & allowed xors & families\\
\midrule
\texttt{DM01} & 4 & 4 & \texttt{3,5,6} & \texttt{NF1}\\
\texttt{DM02} & 4 & 4 & \texttt{2,4,6} & \texttt{NF1}\\
\texttt{DM03} & 6 & 1 & \texttt{3} & \texttt{NF3}\\
\texttt{DM04} & 6 & 1 & \texttt{7} & \texttt{NF2}\\
\texttt{DM05} & 6 & 1 & \texttt{7} & \texttt{NF2}\\
\texttt{DM06} & 6 & 1 & \texttt{7} & \texttt{NF2}\\
\texttt{DM07} & 6 & 2 & \texttt{5,6} & \texttt{NF3}\\
\texttt{DM08} & 6 & 1 & \texttt{7} & \texttt{NF3}\\
\bottomrule
\end{tabular}
\caption{Eight B3 support-action profiles, with allowed selected-label xors and clause families.}
\end{table}

\begin{table}[ht]
\centering
\small
\begin{tabular}{p{0.16\linewidth}p{0.38\linewidth}p{0.34\linewidth}}
\toprule
Certified block & Certified content & Paper consequence\\
\midrule
Intrinsic columns & 8/8 B3 profile columns intrinsic to (S,B) & the column side has no mask-transport residue\\
Defect-6 residue & defect 6 realizes 13 of 35 two-flats; all 35 are F2-completable & the realization cut is integer-magic, not F2-linear\\
Defect-6 residue & the realized 13-set has GL(4,2) stabilizer order 1 & no linear symmetry selects the remaining cut\\
\bottomrule
\end{tabular}
\caption{Post-obstruction B3 status.  The intrinsic-column certificate closes the profile-column side; the defect-6 obstruction identifies the remaining realization cut as integer-magic rather than \(\F_2\)-linear.}
\end{table}

\begin{table}[ht]
\centering
\small
\begin{tabular}{lrl}
\toprule
Family & Clauses & Public rule\\
\midrule
NF1 & 6 & defect 4; singleton xor-error direction\\
NF2 & 3 & defect 6; non-xor pair witness\\
NF3 & 4 & defect 6; non-xor triple witness\\
\bottomrule
\end{tabular}
\caption{NF1/NF2/NF3 clause-template counts.}
\end{table}

\begin{table}[ht]
\centering
\scriptsize
\begin{tabular}{llllll}
\toprule
Clause & Family & Profile & $x$ & $E$ & witness\\
\midrule
\texttt{L20-C01} & NF1 & \texttt{DM01} & 3 & \texttt{3} & \texttt{--}\\
\texttt{L20-C02} & NF1 & \texttt{DM01} & 5 & \texttt{5} & \texttt{--}\\
\texttt{L20-C03} & NF1 & \texttt{DM01} & 6 & \texttt{6} & \texttt{--}\\
\texttt{L20-C04} & NF1 & \texttt{DM02} & 2 & \texttt{2} & \texttt{--}\\
\texttt{L20-C05} & NF1 & \texttt{DM02} & 4 & \texttt{4} & \texttt{--}\\
\texttt{L20-C06} & NF1 & \texttt{DM02} & 6 & \texttt{6} & \texttt{--}\\
\texttt{L20-C07} & NF3 & \texttt{DM03} & 3 & \texttt{3,13,14} & \texttt{14}\\
\texttt{L20-C08} & NF2 & \texttt{DM04} & 7 & \texttt{2,5,7} & \texttt{2,5}\\
\texttt{L20-C09} & NF2 & \texttt{DM05} & 7 & \texttt{1,6,7} & \texttt{1,6}\\
\texttt{L20-C10} & NF2 & \texttt{DM06} & 7 & \texttt{3,4,7} & \texttt{3,4}\\
\texttt{L20-C11} & NF3 & \texttt{DM07} & 5 & \texttt{5,8,13} & \texttt{8}\\
\texttt{L20-C12} & NF3 & \texttt{DM07} & 6 & \texttt{6,8,14} & \texttt{8}\\
\texttt{L20-C13} & NF3 & \texttt{DM08} & 7 & \texttt{7,11,12} & \texttt{11}\\
\bottomrule
\end{tabular}
\caption{The thirteen finite profile-level clause instances.  Empty witnesses are shown as \texttt{--}.}
\end{table}


\FloatBarrier
\section{Certificate Spine}
\label{sec:certificates}

The paper-facing certificate table is maintained in \texttt{paper/SOURCE\_MAP.md}.
Before public export, every displayed number in the manuscript must be tied to
one source artifact and one regression test.  The local test count is deliberately
not hardcoded in the mathematical text; it is a build artifact recorded by the
current regression run.

The key public-safe statement is that the B3 matrix is finite-certified, while
the symbolic derivation of the count \(60\) remains open.  The count-collapse
certificate reduces that debt to the finite endpoint quotient kernel
\[
  \delta=0,4,6 \mapsto 36,16,7
\]
and to the B3 row-degree pattern
\[
  1,1,1,2,2,1,4,3.
\]
The intrinsic column-classification certificate closes the column side by proving that all eight B3 profile columns are
intrinsic to \((S,B)\).  The defect-6 realizability-obstruction certificate does not close the realization side; it proves
that the defect-6 residue is an irreducible integer-magic cut, not a consequence
of \(\F_2\)-completability or linear symmetry.

\begin{certprop}[Reproducibility artifact table]
The manuscript depends on the following certified artifacts.  Exact file names
are kept in \texttt{paper/SOURCE\_MAP.md}; this table uses only reproducibility
handles and mathematical block names.
\begin{center}
\scriptsize
\begin{tabular}{lll}
\toprule
Claim block & Reproducibility handle & Status\\
\midrule
endpoint split & endpoint split certificate & certified\\
defect dichotomy & first-principles defect dichotomy & structural\\
error parity & error-fiber parity bridge & structural\\
B3 matrix & support-action matrix summary & finite\\
B3 intrinsic columns & intrinsic column classification & finite\\
B3 defect-6 residue & defect-6 realizability obstruction & finite\\
endpoint transport & endpoint transport package & finite\\
\(\TK\) freeness & global free-transport theorem & finite\\
B2 package & B2 endpoint boundary package & finite\\
count collapse & count-collapse kernel geometry & finite\\
\bottomrule
\end{tabular}
\end{center}
\end{certprop}

\begin{remark}[Attribution and proof source]
The row/column nim-sum property is positioned against the classical order-four
literature, including Ollerenshaw--Bondi and Ward.  Until a precise external
bibliographic statement is pinned down, the manuscript treats the local Step4
derivation as the proof source.
\end{remark}
\section{Public Companion Package}

The public companion repository follows the curated package structure
\[
\texttt{paper/},\quad \texttt{data/},\quad \texttt{docs/},\quad
\texttt{scripts/},\quad \texttt{tests/},\quad \texttt{results/}.
\]
It is a curated replay and audit package, not a dump of the development
workspace.  Public artifacts are limited to the claims actually made in the
manuscript, and the blocked claims in \texttt{paper/CLAIM\_BOUNDARY.md} remain
blocked.

\paragraph{Blocked claims.}
The public package does not claim public record-level widening, full endpoint
\(23/27\) classification, complete conceptual B2 theory, full B3 non-affine
classification, a symbolic derivation of the B3 \(15/49\) matrix, or a
non-enumerative derivation of the count \(60\).

\paragraph{Release artifacts.}
The package includes curated JSON summaries, source-map and claim-boundary
documents, audit scripts, pytest smoke tests, \texttt{README.md},
\texttt{REPRODUCE.md}, \texttt{CITATION.cff}, and \texttt{LICENSE}.  Its replay
scope is intentionally narrower than the development workspace and traces to
\texttt{paper/SOURCE\_MAP.md}.
\FloatBarrier

\bibliographystyle{plain}
\bibliography{refs}

\end{document}